\chapter*{Chapter 2 glossary}\addcontentsline{toc}{chapter}{Chapter 2 glossary}
\begin{description}
\item[Vertex splint] Complementary strand that aligns adjacent edge fragments at a shared vertex; it does not insert an additional code into the product backbone.
\item[Nick] Discontinuity in one strand's backbone within an otherwise paired structure. Pairing across the junction does not itself seal the nick.
\item[Affinity selection] Recognition and retention using a binding interaction. Here a complementary probe recognizes a vertex code; biotin attaches the probe to a recoverable support.
\item[Graduated PCR] A set of primer-dependent amplifications whose product lengths can support position inference. The lane projection does not retain joint identity across a mixture.
\item[Soundness and completeness] Correctness of accepted witnesses versus recovery and detection of existing witnesses. A physical implementation must consider both.
\item[Support and multiplicity] Which route types are present versus how many copies of each exist. Neither is determined solely by total reagent molecules.
\end{description}
